\documentclass[11pt,letterpaper]{article}
\usepackage[margin=1in]{geometry}
\usepackage{hyperref}
\usepackage{parskip}
\usepackage{amsmath,amssymb}
\usepackage{titlesec}
\usepackage{booktabs}

\title{\textbf{Seeing the Bayes Net in the Transformer: A Tutorial}}
\author{Greg Coppola\\coppola.ai\\greg@coppola.ai}
\date{2026}

\begin{document}
\maketitle

\begin{abstract}
The shannon paper (arXiv:2603.17063) proved formally that a sigmoid transformer is a
Bayesian network --- not approximately, not analogically, but architecturally, for any
weights, Lean-verified. This tutorial asks the follow-up question: can you actually
\emph{see} the Bayes net when you look at a real trained transformer? Can you point at
specific weights, heads, and layers and say --- that is the AND gate, that is the OR
gate, that is the factor potential?

We argue the answer is yes. We walk through five correspondences between the transformer
architecture and belief propagation on a factor graph: the residual stream as the belief
state, attention as the gather step (AND), the FFN as the update step (OR), layers as
rounds of BP, and the scaling numbers as a Bayes net inventory. We address the
dimensionality question --- why $D_{ff} \approx 3$--$4 \times D$ --- via superposition,
$k$-ary neighbors, and activation precision. We survey the mechanistic interpretability
literature and show that every major finding (induction heads, key-value memories, ROME,
SAE features, attribution graphs) maps naturally onto a component of the BP account.
Five independent research programs have converged on this answer. This tutorial explains
what the answer means.
\end{abstract}

\tableofcontents
\newpage

\section{Introduction}
% TODO: full introduction prose
% Key points: the shannon result, the tutorial question, preview of five correspondences,
% the convergence story, tone setting.
\textit{[Section to be written. See index.tex for content summary.]}

\section{Background: Factor Graphs and Belief Propagation}
% TODO: self-contained BP intro
% Covers: variable nodes, factor nodes, joint distribution, gather-update algorithm,
% updateBelief formula, convergence on trees, loopy approximation.
\textit{[Section to be written. See index.tex for content summary.]}

\section{Background: The Transformer}
% TODO: minimal transformer account in BP framing
% Covers: residual stream, attention, FFN, five dimensions (W, D, H, D_ff, L),
% routing view (Elhage 2021).
\textit{[Section to be written. See index.tex for content summary.]}

\section{The Log-Odds Algebra}
% TODO: Turing-Good-Pearl unification
% Covers: logit/sigmoid isomorphism, the algebra on (0,1), three independent derivations,
% why sigmoid is the right activation, relation to Boolean logic.
\textit{[Section to be written. See index.tex for content summary.]}

\section{The Five Correspondences}

\subsection{The Residual Stream as the Belief State}
% TODO
\textit{[Section to be written. See index.tex for content summary.]}

\subsection{Attention as the Gather Step (AND)}
% TODO
\textit{[Section to be written. See index.tex for content summary.]}

\subsection{The FFN as the Update Step (OR)}
% TODO
\textit{[Section to be written. See index.tex for content summary.]}

\subsection{Layers as Rounds of BP}
% TODO
\textit{[Section to be written. See index.tex for content summary.]}

\subsection{The Scaling Numbers as a Bayes Net Inventory}
% TODO
\textit{[Section to be written. See index.tex for content summary.]}

\section{The Dimensionality Question: $D_{ff}$, Superposition, and Multiple OR Gates}
% TODO: the D_ff/D ratio, why it is 3-4x, three mechanisms
\textit{[Section to be written. See index.tex for content summary.]}

\section{Evidence from Mechanistic Interpretability}
% TODO: each finding + BP interpretation
\textit{[Section to be written. See index.tex for content summary.]}

\section{The Three Softmaxes}
% TODO: routing vs inference vs generation
\textit{[Section to be written. See index.tex for content summary.]}

\section{Grounded vs.\ Ungrounded Transformers}
% TODO: grounding, hallucination in BP terms, RAG as partial grounding
\textit{[Section to be written. See index.tex for content summary.]}

\section{Open Questions}
% TODO: diffuse attention, weight inspection experiment, grounding gap
\textit{[Section to be written. See index.tex for content summary.]}

\section{Related Work}
% TODO: five clusters, convergence table
\textit{[Section to be written. See index.tex for content summary.]}

\section{Conclusion}
% TODO
\textit{[Section to be written. See index.tex for content summary.]}

\bibliographystyle{plain}
\bibliography{refs}

\end{document}